% A report fragment; requires amsmath, amssymb, booktabs and hyperref.
% New proof for the coefficient-five Q / no-push MP / cached-only Shift composition.
\section{The PA194QH5 construction}
\label{sec:qh5}

PA194QH5 combines three established source features that the pinned corpus
has not combined at coefficient five: premise-exchanged quantifier
distribution, successful modus ponens without a dummy push, and a Shift
that tests only the cached formula. Its formula constructors use coefficient
five. The components are prior work. The new result is this exact
composition, the complete literal binary transition table, and its checked
source, physical-runtime, and initialization bindings.

There are 194 working states, numbered $0,\ldots,193$, and a separate halt
label 194. The start is state zero on an all-zero two-way infinite tape.
Every transition writes one of $0,1$ and moves exactly one cell left or
right. There is one tape, no stay-put move, and no preloaded input.
In native numbering the start is 74 and the halt label is 254.
The table is built from 135 fixed controls and 67 reader controls, with
eight simultaneous history folds:
\[
                       135+67-8=194.
\]
The builder checks all 32 combinations of history and current bits for
these folds on the final simultaneous table. All working states are
reachable in the directed action graph from the designated start; this
finite graph fact does not assert that the blank computation visits them all.

\subsection{The exact change and its source meaning}

Relative to the native PA194QH6 table, exactly one transition changes:
\[
 \delta_6(190,1)=(1,R,330),\qquad
 \delta_5(190,1)=(1,R,99).
\]
The change shortens the common formula-construction body from six
increments to five. Write $i$ for the one-increment source leaf and
$ii$ for its two-increment reader. The old and new binary source trees
are respectively $((i,ii),(i,ii))$ and $((i,ii),ii)$.
The changed branch is in a recurring arithmetic operation, not an unused
transition, a relabeling, or a change of tape direction.

The complete source expands to 3,714 primitive sites and 4,815 outcome
contracts. The independently resolved successor graph is part of the
release. There are precisely 61 occurrences of the changed body, and
precisely 61 formula constructors; the body occurs at relative path
$1110$ inside each constructor. For every occurrence the builder checks
its entire 21-site unpair-and-pair prefix, decrement guard, five-increment
body, and back edge. Every one of the other 3,409 primitive sites has the
same operation and register argument as in QH6. All their resolved
successors agree under an explicit address relocation. No edge from
outside a body enters its interior. Source entry 25 and its binary
address are unchanged.

These finite facts establish where the arithmetic change occurs, but its
meaning requires an unbounded induction. In the final constructor loop,
with $T=r_1$ and $Y=r_7$, a positive iteration sends
$(T,Y)$ to $(T+5,Y-1)$. After $j\leq Y_0$ iterations,
\[
                   T=T_0+5j,\qquad Y=Y_0-j.
\]
The loop therefore ends with $(T_0+5Y_0,0)$ in $7Y_0+1$ primitive
instructions. The preceding unpair-and-pair phase computes the ordinary
Cantor pair and clears the scratch register. Consequently, with
$[a;w]=p(a,w)$ and scratch $X=Y=0$, the macro has the all-input contract
\[
 (T,W,X,Y)=(t,[a;w],0,0)
 \longmapsto (5p(a,t),w,0,0).
\]
The displayed coordinates are constructor registers, not the command
store coordinates below. The full unpair, pair, and multiplication
induction is preserved in \texttt{references/logical/generic\_cons\_proof.tex}.
The source binding checks that the quoted prefix is exactly the prefix
used here. Stack pairing and scheduler pairing remain ordinary $p$;
only formula construction uses the factor five. Modus-ponens matching
passes through the same constructor, so it uses the same factor.

\subsection{The numerical calculus}

Let $p(a,b)=(a+b)(a+b+1)/2+b$. Define
\[
 I(A,B)=5p(A,B),\quad M(a,b)=5p(a,b)+1,\quad
 U(a,B)=5p(B,a)+2,\quad E(a,b)=5p(a,b)+3.
\]
The order of the two arguments in $U$ is deliberate. In displays write
$A\to B$ for $I(A,B)$, $N(B)=I(B,1)$,
$A\wedge B=N(I(A,N(B)))$, and $\exists_a B=N(U(a,N(B)))$.
Zero denotes truth and one is the designated contradiction $M(0,0)$.
The numerical alias $I(0,0)=0$ is retained throughout the proof.

The six ordinary axiom families are
\begin{align*}
 h_1(D,H)&=(\exists_D U(D,H))\to H,\\
 h_2(D,A,H)&=U(D,U(A,H))\to U(A,U(D,H)),\\
 h_3(D,A,H)&=M(H,D)\to(E(H,A)\to M(A,D)),\\
 h_4(D,A,H)&=E(H,D)\to(M(A,H)\to M(A,D)),\\
 h_5(D,H)&=\exists_D E(H,D),\\
 h'_6(D,A,H)&=U(D,A\to H)\to(U(D,A)\to U(D,H)).
\end{align*}
There is also the restricted implication family
\[
 \rho_8(D,A,H,u)=((H\to u)\to D)\to((D\to H)\to(A\to H)),
 \qquad u\leq D,
\]
and the diagonal family $\Gamma(B)=I(B,U(B,B))$. Finally put
\begin{align*}
 V(a,b,c)&=M(a,b)\to U(c,M(a,b)),\\
 P(a,b,c)&=c\to M(c,b),\\
 S(a,b,c)&=U\bigl(b,(M(a,b)\to P(a,b,c))
                         \wedge(M(c,b)\to c)\bigr),\\
 Q(a,b,c)&=U(c,S(a,b,c)\to N(P(a,b,c))),\\
 G(a,b,c)&=N(V(a,b,c)\to Q(a,b,c)).
\end{align*}
The gate family consists of $G(a,b,c)$ for $c\ne a,b$. The only inference
rules are modus ponens and theorem generalization $B\mapsto U(i,B)$.
There are no undischarged hypotheses; formulas may have free labels.
Call this calculus $\mathcal D_5$. Failed gate emissions return zero,
which is already a theorem since $\rho_8(0,0,0,0)=0$.

The ordinary coefficient-five separated calculus $\mathcal C_5$ differs
only in having
$h_6=U(D,A)\to(U(D,A\to H)\to U(D,H))$.
The unchanged restricted implication basis derives $K,S$ and the fixed
exchange template
\[
  (P\to(Q\to R))\to(Q\to(P\to R)).
\]
A 61-line $K,S,\mathrm{MP}$ certificate is supplied and independently
replayed. Applying exchange in either direction replaces one version of
$h_6$ by the other. Induction on finite derivations gives two
primitive-recursive, conclusion-preserving proof translations whose
correctness is provable in PA. Thus
\[
 \mathrm{PA}\vdash
 \bigl(\exists d\,\operatorname{Prf}_{\mathcal D_5}(d,1)
       \leftrightarrow\exists d\,\operatorname{Prf}_{\mathcal C_5}(d,1)\bigr).
\]
Here and below $\operatorname{Prf}$ denotes the relevant primitive-recursive proof
checking relation. No unrestricted formula substitution into the guarded
gate is used in the exchange argument.

\subsection{Calibration to PA, with all numerical edge cases}

The inherited mathematical dependency is the PA calibration of the
coefficient-four separated diagonal calculus $\mathcal C_4$. Its complete
supplied calibration arguments are retained in the reference directory:
original PA195, diagonal-calculus replacement, and separated-head
calibration. They supply primitive-recursive refutation translations
between $\mathcal C_4$ and PA, with correctness proved by finite syntactic
inductions in PA. These are ordinary mathematical arguments; their
presence is not a claim of a newly completed Lean or Rocq endpoint.

For completeness at coefficient five, define $F(0)=0$ and, for $n>0$,
\[
 n=4p(a,b)+r\ (r<4)\quad\Longrightarrow\quad
 F(n)=5p(F(a),F(b))+r.
\]
Both recursive arguments are smaller than $n$. This course-of-values
recursion is primitive recursive: a primitive recursion can retain all
previous values and retrieve the required entries. Define $J(0)=0$ and,
for positive inputs,
\[
 J(5p(a,b)+r)=
 \begin{cases}4p(J(a),J(b))+r,&r<4,\\0,&r=4.\end{cases}
\]
PA proves $J(F(n))=n$ by induction; hence $F$ is injective and preserves
the disequality guards. It fixes zero and one and commutes with every
constructor. The same $F$ must act on formula codes and variable labels.
In particular $F(\Gamma_4(B))=\Gamma_5(F(B))$, including the
self-indexed binder. This is not a recoding that holds labels fixed.

Order is not preserved: $40<44$, but $F(40)=75>55=F(44)$.
A restricted propositional axiom therefore cannot be translated by
silently assuming $F(u)\leq F(D)$. The inherited finite
\texttt{rho8All} derivation supplies unrestricted $\rho_8$ from the
restricted basis. Its only order operations are reflexivity and the
constructor bounds
\[
                        a,b\leq p(a,b)\leq5p(a,b).
\]
The finite template is retained verbatim in
\texttt{references/lean/M6/Logic/RestrictedImplication.lean}.
At a recoded restricted leaf, insert this unrestricted template;
its own restricted leaves satisfy their actual target inequalities.
Recursing on the source proof now transfers every $\mathcal C_4$
refutation to a $\mathcal C_5$ refutation. This argument uses neither
monotonicity of $F$ nor preservation of subtraction or stack pairing.

For the reverse reduction, define a total \emph{syntactic} arithmetic
translation $T$. Send zero and every unused-residue input $n\equiv4$
modulo five to $0=0$. At a positive input $n=5p(a,b)+r$, the four
used tags have translations
\[
\begin{array}{c|l}
 r& T(n)\\\hline
 0&T(a)\to T(b)\\
 1&\operatorname{Bit}(x_a,x_b)\\
 2&\forall x_b\,T(a)\\
 3&x_a=x_b.
\end{array}
\]
Use capture-avoiding binders for the internal bit predicate. This is a
primitive-recursive formula-code constructor, not an arithmetic truth
predicate. Every source variable label occurring in $T(B)$ is strictly
below $B$; the proof is induction on the strictly smaller decoded
components, with zero and unused residue interpreted as closed truth.
Thus $x_B$ is absent from $T(B)$ and $\Gamma(B)$ has a uniform PA proof
by vacuity. The constructor equivalence
$T(I(A,B))\leftrightarrow(T(A)\to T(B))$ is literal except at
$A=B=0$, where the fixed truth/implies-truth equivalence supplies the
proof. No injective syntax assumption is made at code zero.

The bit predicate may be written
\[
 \operatorname{Bit}(s,z)\equiv
 \exists u,v\,[s=(2u+1)2^z+v\ \wedge\ v<2^z].
\]
PA proves $\operatorname{Bit}(s,z)\to z<s$, so it refutes $T(1)$.
For every arithmetic formula $\theta(z,\vec u)$, substitution into a
fixed PA induction proof gives a primitive-recursive supplier of a proof
of finite binary comprehension:
\[
 \forall n\,\exists s<2^n\,\forall z<n\,
       (\operatorname{Bit}(s,z)\leftrightarrow\theta(z,\vec u)).
\]
The induction step retains the old bit set when $\neg\theta(n,\vec u)$
and adds $2^n$ otherwise. This syntactic supplier uses the corresponding
full PA induction instance. It neither decides $\theta$ nor assumes
reflection for PA.

For the guarded gate let $\phi=T(c)$. If the labels $a,b$ are distinct,
use the finite-comprehension bound $x_a+x_b+1$ and matrix
$(\operatorname{Bit}(x_a,z)\vee z=x_b)\wedge\phi[z/x_b]$.
If $a=b$, use bound $x_b+1$ and matrix
$(z=x_b)\wedge\phi[z/x_b]$. In the coincident-label case the first
membership antecedent under the binder is $\operatorname{Bit}(z,z)$,
which PA refutes. The chosen set satisfies $S$ and the point condition
$P$, hence refutes $Q$. The guard $c\ne a,b$ supplies membership
vacuity $V$, and freshness of label $c$ in $T(c)$ prevents the set
witness from changing $\phi$. Thus PA proves the interpreted gate.
This explicitly treats the coincident-label case rather than confusing
the bound occurrence of $a$ with a free value.

All six ordinary heads, the exchanged sixth head, and the restricted
implication head have uniform first-order or propositional proof
templates. Annotations transform MP by the constructor equivalence and
generalization by the PA generalization rule. Recursion on a finite
$\mathcal D_5$ derivation therefore converts a proof of one to a PA
refutation. Combining the directions gives the ordinary syntactic theorem
\[
 \mathrm{PA}\vdash
 \left(\exists d\,\operatorname{Prf}_{\mathcal D_5}(d,1)
  \leftrightarrow\exists p\,\operatorname{Prf}_{\mathrm{PA}}(p,\ulcorner0=1\urcorner)\right).
 \tag{C5}
\]
The universal statement is about proof codes in PA, including nonstandard
codes in a model, rather than a separate external verification for each
standard proof. The general radix proof is retained in
\texttt{references/logical/radix\_calibration.tex}; its older controller
section is historical context and is superseded here by the next section.

\subsection{Command semantics and an all-store proof compiler}

A command store is $(t,w,d,h)$, with top formula $t$, ordinary Cantor
stack $w$, offset $d$, and cached formula $h$. Its two non-axiom commands
are, with $b=d+x$,
\[
 \mathsf S(t,w,d,h)=
 \begin{cases}\mathsf{halt},&h=1,\\
 (h,\operatorname{tl}w,d,\operatorname{hd}w),&h\ne1,
 \end{cases}
\]
\[
 \mathsf{MP}_x(t,w,d,h)=
 \begin{cases}(b,w,b,0),&I(t,b)=h,\\
 (0,w,b,0),&I(t,b)\ne h.
 \end{cases}
\]
Successful MP does not insert a dummy stack entry. With $D=d+x$,
$H=h+y$, $A=c$, the axiom command has effect
\[
 \mathsf{AX}_{c,x,y}(t,w,d,h)=
 (G(H,D,A),[q_1,\ldots,q_9,t;w],0,0),
\]
where a failed guard returns zero and
\[
 (q_1,\ldots,q_9)=
 (\Gamma(A),h'_6,h_5,h_4,h_3,h_2,h_1,
                        \rho_8(D,A,H,d),U(d,h)).
\]
All $q_1,\ldots,q_8$ and the emitted gate have residue zero modulo five;
$q_9$ has residue two. None is one. The complete free-constructor emitter
check verifies all 190 macro tokens and the entire output stack and top
for arbitrary symbolic inputs.

Let $A_0=\mathsf{AX}_{0,0,0}$. Define
\[
 \mathsf{Push}=A_0;\mathsf S^9;\mathsf{MP}_0,\quad
 \mathsf{Pop}=A_0;\mathsf S^{11},\quad
 \mathsf{Discard}=\mathsf{MP}_0,\quad
 \mathsf{Offset}_i=\mathsf{MP}_i.
\]
Direct calculation proves, for all displayed natural inputs,
\begin{align*}
 \mathsf{Push}(t,w,d,h)&=(0,[t;w],0,0),\\
 \mathsf{Pop}(t,[a;w],d,h)&=
    \begin{cases}\mathsf{halt},&t=1,\\(t,w,0,a),&t\ne1,\end{cases}\\
 \mathsf{Discard}(t,w,0,h)&=(0,w,0,0),\\
 \mathsf{Offset}_i(0,w,0,0)&=(0,w,i,0).
\end{align*}
Push's comparison fails by the implication/universal residue mismatch;
Discard's two branches coincide at proposed conclusion zero. The first
ten Shifts after AX test only the cleared cache and generated entries;
the eleventh tests exactly the saved top. The identities include
$[0;0]=0$ and the total projections of zero, both equal to zero.

From $(t,w,0,0)$, the head selector for $1\leq j\leq7$ is
\[
 \mathsf{Discard};\mathsf{AX}_{A,D,H};\mathsf S^{j+1};
                    \mathsf{Pop}^{9-j}.
\]
It returns $q_j$, preserves $w$, and clears caches. The gate selector is
$\mathsf{Discard};\mathsf{AX}_{c,b,a};\mathsf{Pop}^{10}$.
For a restricted implication leaf with $u\leq D$, use
$\mathsf{Discard};\mathsf{Offset}_u;
\mathsf{AX}_{A,D-u,H};\mathsf S^9;\mathsf{Pop}$.
This subtraction is on actual valid target parameters, after the
unrestricted-template expansion described above. Generalization uses
\[
 \mathsf{Gen}_i=\mathsf{Push};\mathsf{Offset}_i;\mathsf S;
              \mathsf{AX}_{i,0,0};\mathsf S^{10},
\]
and sends $(a,w,0,0)$ to $(U(i,a),w,0,0)$.

For MP, the suffix $\mathsf S;\mathsf S;\mathsf{MP}_b;\mathsf{Pop}$
on $(0,[a,I(a,b),0;w],0,0)$ returns $(b,w,0,0)$ if $a,b\ne1$,
and halts if either equals one. The protected zero below the premises
replaces the old controller's dummy push. At a proof-tree MP node with
major proof $p_1$ and minor proof $p_0$, take
\[
 W_p=\mathsf{Discard};\mathsf{Push};W_{p_1};\mathsf{Push};
 W_{p_0};\mathsf{Push};\mathsf S;\mathsf S;\mathsf{MP}_b;\mathsf{Pop}.
\]
Induction on a proof establishes that from every $(t,w,0,0)$ its word
either halts or returns its conclusion, preserving $w$ and clearing the
caches. Thus a refutation yields $A_0;W_p;\mathsf{Pop}$, which halts
from \emph{every incoming store}. A final plain Shift would fail:
$\mathsf S(1,0,0,0)=(0,0,0,0)$. The initial AX is necessary to clear
arbitrary carried caches. The separate symbolic command checker verifies
the exact identities, including their zero aliases and both MP branches.

\subsection{Fair scheduling, proof annotations and the literal machine}

The inherited scheduler retains the store between trials, so the all-store
property is essential. Its ordinary Cantor encoding is $e(\epsilon)=0$
and $e(c::w)=1+b_c(e(w))$, where
\[
 b_{\mathsf S}(z)=p(0,z),\quad
 b_{\mathsf{MP}_x}(z)=p(1,p(x,z)),\quad
 b_{\mathsf{AX}_{c,x,y}}(z)=p(c+2,p(y,p(x,z))).
\]
For an all-store halting suffix of continuation code $z$, prefixing
$\mathsf{AX}_{k,0,0}$ gives the unbounded family of candidate codes
$p(k+2,p(0,p(0,z)))$. Fresh candidates begin at two. Every positive
continuation decreases strictly; the register arithmetic has
primitive-recursive termination bounds. PA verifies by finite induction
that a member of this family is reached unless an earlier halt occurs.
Thus a $\mathcal D_5$ refutation yields a finite command halt.

For soundness annotate $t,h$ and every canonical stack entry with a
$\mathcal D_5$ proof. Positive Cantor-stack tails decrease strictly, so
the list and its annotations have primitive-recursive codes. Virtual
zero entries use the fixed proof of zero. AX emits its axioms, the old
top, and a generalization of the old cache. Successful MP combines
annotated premises; failure returns zero; Shift moves annotations and
halts only at a proof of one. Finite induction proves that every command
halt yields a $\mathcal D_5$ refutation, hence by (C5) a PA refutation.

The construction checker verifies all 4,815 primitive source outcomes
for arbitrary nine primary registers and both decrement branches, using
1,576,144 symbolic operations with maximum case 499. A separately
implemented checker derives the four generic register-list lemmas from
the literal transition rows and replays the same complete source using
two independent successor resolvers. Its domain has arbitrary natural
registers, arbitrary nonempty finite spare-register list, and affine
header gap $48+n-|p|$ for every $n\geq0$. The maximum source address
length is 36. These are universal symbolic certificates, not finite
sampling of register values. Their interpretation is justified by
induction on unary runs and register lists.

Exactly 58,088 literal transitions from blank reach the clean source
entry at native state 84 and head $-2$, logical primary registers
$(2,0,0,0,0,0,0,0,0)$, and 26 unit spare registers. Both checkers compare
the entire finite-support tape and its blank exterior. No initialization
is supplied outside the counted states.

A finite source halt therefore gives a literal blank-tape halt by
concatenation of positive-duration physical simulations. Conversely,
given a putative first literal halting time $T$, if the command source
never halted, concatenate $T+1$ positive completed source segments after
startup. Their deterministic literal execution extends beyond $T$
without halting, a contradiction. This finite argument is formalizable
in PA and avoids an inverse simulation of arbitrary intermediate tapes.
Combining it with (C5), the compiler and the annotation invariant gives
\[
 \boxed{\mathrm{PA}\vdash
 \bigl(\neg\operatorname{Halt}_{\rm blank}(\mathrm{PA194QH5})
                  \leftrightarrow\operatorname{Con}(\mathrm{PA})\bigr).}
\]
This is an ordinary mathematical theorem using the explicitly identified
inherited PA calibration. Executable checks verify the new finite
bindings and universal symbolic source contracts. No end-to-end Lean,
Rocq or iPRA theorem for QH5 is claimed. Actual nonhalting follows if PA
is consistent; no simulation can establish that premise.

\subsection{Materiality and limits}

The factor five changes every formula-construction use, including the
axiom emitter and MP comparison. For example $I_5(1,1)=20$ whereas
$I_6(1,1)=24$. This is a change in recurring arithmetic source semantics.
Relative to earlier coefficient-five Q5, QH5 also changes both command
policies; relative to N5, it changes the distribution axiom and the Shift
test. It is a new composition of established ideas, not a new component
or a claimed lower bound. The exact action-graph census provides the
separate finite test against state renaming, tape reflection, and symbol
exchange. Formula recoding itself is not claimed to conjugate addition,
stacks, stores or individual trajectories; it transports proofs, which
are then recompiled into new target words.
